\documentclass[aspectratio=43,10pt]{beamer}
\usetheme{Montpellier}
\usepackage[utf8]{inputenc}
\usepackage[T1]{fontenc}
\usepackage[brazil]{babel}
\usepackage{amsmath,amssymb,booktabs,listings,url,tikz,fontawesome5,graphicx}
\usetikzlibrary{arrows.meta,positioning,shapes.geometric}

% ---------------------------------------------------------------------------
% Paleta e estilos TikZ
% ---------------------------------------------------------------------------
\definecolor{treeblue}{RGB}{34,92,135}
\definecolor{treegreen}{RGB}{55,122,86}
\definecolor{treeorange}{RGB}{207,119,37}
\definecolor{treered}{RGB}{176,60,60}

\tikzset{
  node/.style={circle,draw=treeblue,fill=blue!9,thick,minimum size=6.5mm,inner sep=0pt,font=\small},
  key/.style={node,draw=treeorange,fill=orange!17},
  leaf/.style={node,draw=treegreen,fill=green!13},
  edge/.style={draw=treeblue,thick,-{Stealth[length=2mm]}},
  plain/.style={draw=treeblue,thick},
  decision/.style={diamond,draw=treeblue,fill=blue!9,thick,align=center,aspect=2,font=\scriptsize,inner sep=1.5pt},
  outcome/.style={rounded corners,draw=treegreen,fill=green!13,thick,align=center,font=\scriptsize,minimum height=6mm},
  folder/.style={draw=treeblue,rounded corners,fill=blue!9,thick,align=center,font=\scriptsize,inner sep=3pt},
  filenode/.style={draw=treegreen,rounded corners,fill=green!12,thick,align=center,font=\scriptsize,inner sep=2.5pt},
  % Espaçamento por nível: a distância entre irmãos do nível 2 é
  % deliberadamente menor que a do nível 1, para que "primos" de subárvores
  % vizinhas nunca fiquem colineares e se sobreponham.
  arv/.style={
    level distance=11mm,
    level 1/.style={sibling distance=20mm},
    level 2/.style={sibling distance=9mm},
    edge from parent/.style={draw=treeblue,thick},
  },
}

\lstset{basicstyle=\ttfamily\scriptsize,keywordstyle=\bfseries\color{treeblue},commentstyle=\itshape\color{treegreen},showstringspaces=false,frame=single,breaklines=true,columns=fullflexible}

\title[Estruturas de Dados: Árvores]{Estruturas de Dados: Árvores}
\author[Elton Sarmanho]{Professor: Elton Sarmanho\inst{1}\\\small \faEnvelope\,\href{mailto:eltonss@ufpa.br}{eltonss@ufpa.br}}
\institute[UFPA]{\inst{1}\faUniversity\, Faculdade de Sistemas de Informação -- UFPA/CUTINS}
\date{\today}
\setbeamertemplate{navigation symbols}{}
\setbeamertemplate{footline}{\hfill\insertshorttitle\hspace{1em}\insertframenumber/\inserttotalframenumber\hspace{1em}\vspace{.5em}}
\setlength{\columnsep}{6pt}

% ---------------------------------------------------------------------------
% Macros de diagramas reutilizáveis
% ---------------------------------------------------------------------------

% Árvore de referência A-B-C-D-E-F-G, usada em vários slides conceituais.
\newcommand{\ABC}{%
\begin{tikzpicture}[arv,baseline=(current bounding box.center)]
\node[node]{A}
  child{node[node]{B} child{node[leaf]{D}} child{node[leaf]{E}}}
  child{node[node]{C} child{node[leaf]{F}} child{node[leaf]{G}}};
\end{tikzpicture}}

% Mesma árvore, com nível anotado entre parênteses (usa espaçamento maior
% porque os rótulos "A (0)" etc. são mais largos que uma letra isolada).
\newcommand{\ABCNiveis}{%
\begin{tikzpicture}[
    level distance=13mm,
    level 1/.style={sibling distance=26mm},
    level 2/.style={sibling distance=13mm},
    edge from parent/.style={draw=treeblue,thick},
    baseline=(current bounding box.center)]
\node[key]{A\,(0)}
  child{node[node]{B\,(1)} child{node[leaf]{D\,(2)}} child{node[leaf]{E\,(2)}}}
  child{node[node]{C\,(1)} child{node[leaf]{F\,(2)}} child{node[leaf]{G\,(2)}}};
\end{tikzpicture}}

% Árvore binária de busca 50/30/70/20/40/60/80.
\newcommand{\BST}{%
\begin{tikzpicture}[arv,baseline=(current bounding box.center)]
\node[key]{50}
  child{node[node]{30} child{node[leaf]{20}} child{node[leaf]{40}}}
  child{node[node]{70} child{node[leaf]{60}} child{node[leaf]{80}}};
\end{tikzpicture}}

% Pequena árvore de diretórios que ilustra hierarquia de arquivos.
\newcommand{\FileTreeDiagram}{%
\begin{tikzpicture}[
    level distance=10mm,
    level 1/.style={sibling distance=27mm},
    edge from parent/.style={draw=treeblue,thick},
    baseline=(current bounding box.center)]
\node[folder]{\faFolderOpen\ /}
  child{node[folder]{\faFolder\ Aulas} child{node[filenode]{\faFile\ Slides}}}
  child{node[folder]{\faFolder\ Dados} child{node[filenode]{\faFile\ Notas}}};
\end{tikzpicture}}

% Remoção em ABB, caso "nó folha": remover 3 de {4,2,6,1,3,7}.
\newcommand{\DeleteLeafDiagram}{%
\begin{tikzpicture}[x=6.3mm,y=6.3mm,baseline=(current bounding box.center)]
\node[font=\small] at (0,1.0) {Antes: remover 3};
\node[node] (r1) at (0,0) {4};
\node[node] (l1) at (-1,-1) {2};
\node[node] (q1) at (1,-1) {6};
\node[leaf] (a1) at (-1.5,-2) {1};
\node[leaf,draw=treeorange,fill=orange!18] (x1) at (-.5,-2) {3};
\node[leaf] (z1) at (1.5,-2) {7};
\draw[plain] (r1)--(l1) (r1)--(q1) (l1)--(a1) (l1)--(x1) (q1)--(z1);
\node[font=\small] at (5,1.0) {Depois};
\node[node] (r2) at (5,0) {4};
\node[node] (l2) at (4,-1) {2};
\node[node] (q2) at (6,-1) {6};
\node[leaf] (a2) at (3.5,-2) {1};
\node[leaf] (z2) at (6.5,-2) {7};
\draw[plain] (r2)--(l2) (r2)--(q2) (l2)--(a2) (q2)--(z2);
\end{tikzpicture}}

% Remoção em ABB, caso "um filho": remover 20 de {10,5,20,30}.
\newcommand{\DeleteOneChildDiagram}{%
\begin{tikzpicture}[node distance=23mm,baseline=(current bounding box.center)]
\node[font=\small] (a) {Antes: remover 20};
\node[font=\small,right=44mm of a] (b) {Depois};
\node[node,below=7mm of a] (r1) {10};
\node[leaf,below left=11mm and 13mm of r1] (l1) {5};
\node[key,below right=11mm and 13mm of r1] (x1) {20};
\node[leaf,below=11mm of x1] (y1) {30};
\draw[plain] (r1)--(l1) (r1)--(x1) (x1)--(y1);
\node[node,below=7mm of b] (r2) {10};
\node[leaf,below left=11mm and 13mm of r2] (l2) {5};
\node[leaf,below right=11mm and 13mm of r2] (y2) {30};
\draw[plain] (r2)--(l2) (r2)--(y2);
\end{tikzpicture}}

% Remoção em ABB, caso "dois filhos": remover 15, sucessor 20.
\newcommand{\DeleteTwoChildrenDiagram}{%
\begin{tikzpicture}[baseline=(current bounding box.center)]
\node[key] (r) {15};
\node[node,below left=11mm and 13mm of r] (l) {10};
\node[node,below right=11mm and 13mm of r] (q) {25};
\node[leaf,below left=11mm and 8mm of q,draw=treeorange,fill=orange!18] (s1) {20};
\node[leaf,below right=11mm and 8mm of q] (s2) {30};
\draw[plain] (r)--(l) (r)--(q) (q)--(s1) (q)--(s2);
\end{tikzpicture}}

% Rotação simples à esquerda (caso RR): inserir 12 sob 6-8.
\newcommand{\RotationLeftDiagram}{%
\begin{tikzpicture}[x=6.3mm,y=6.3mm,baseline=(current bounding box.center)]
\node[font=\scriptsize] at (0,1.0) {Antes};
\node[key] (a) at (0,0) {6};
\node[node] (b) at (1,-1) {8};
\node[leaf] (c) at (2,-2) {12};
\draw[plain] (a)--(b)--(c);
\node[font=\large] at (3.1,-.7) {\faArrowRight};
\node[font=\scriptsize] at (5.8,1.0) {Depois};
\node[key] (r) at (5.8,0) {8};
\node[leaf] (l) at (4.8,-1) {6};
\node[leaf] (q) at (6.8,-1) {12};
\draw[plain] (r)--(l) (r)--(q);
\end{tikzpicture}}

% Rotação dupla direita-esquerda (caso RL): inserir 7 sob 6-8.
\newcommand{\RotationRLDiagram}{%
\begin{tikzpicture}[x=6.3mm,y=6.3mm,baseline=(current bounding box.center)]
\node[font=\scriptsize] at (0,1.0) {Antes};
\node[key] (a) at (0,0) {6};
\node[node] (b) at (1,-1) {8};
\node[leaf] (c) at (.2,-2) {7};
\draw[plain] (a)--(b) (b)--(c);
\node[font=\large] at (3.1,-.7) {\faArrowRight};
\node[font=\scriptsize] at (5.8,1.0) {Depois};
\node[key] (r) at (5.8,0) {7};
\node[leaf] (l) at (4.8,-1) {6};
\node[leaf] (q) at (6.8,-1) {8};
\draw[plain] (r)--(l) (r)--(q);
\end{tikzpicture}}

% Árvore de decisão de crédito (renda / histórico). Diamantes com largura
% fixa e rótulos em duas linhas para caber no espaço da coluna; as etiquetas
% sim/não ficam em selos com fundo branco sobre a aresta, sem tocar os nós.
\newcommand{\DecisionTreeDiagram}{%
\resizebox{.94\linewidth}{!}{%
\begin{tikzpicture}[
    x=1mm,y=1mm,
    decision/.style={diamond,draw=treeblue,thick,fill=blue!6,align=center,
      font=\fontsize{5.6}{6.4}\selectfont,inner sep=1pt,minimum width=20mm,minimum height=11.5mm},
    outok/.style={rounded corners=1.6pt,draw=treegreen,thick,fill=green!14,align=center,
      font=\fontsize{5.6}{6.4}\selectfont,minimum width=16mm,minimum height=6.5mm,inner sep=1.5pt},
    outwarn/.style={rounded corners=1.6pt,draw=treeorange,thick,fill=orange!16,align=center,
      font=\fontsize{5.6}{6.4}\selectfont,minimum width=16mm,minimum height=6.5mm,inner sep=1.5pt},
    outno/.style={rounded corners=1.6pt,draw=treered,thick,fill=red!7,align=center,
      font=\fontsize{5.6}{6.4}\selectfont,minimum width=16mm,minimum height=6.5mm,inner sep=1.5pt},
    lbl/.style={font=\fontsize{5}{5}\selectfont,fill=white,inner sep=1.2pt,rounded corners=1pt},
    baseline=(current bounding box.center)]
\node[decision] (income)  at (0,0)     {Renda mensal\\$>$ R\$\,3\,mil?};
\node[decision] (history) at (-12,-19) {Histórico de\\crédito bom?};
\node[outno]    (reject)  at (14,-19)  {\faTimesCircle\ Não aprovar};
\node[outok]    (approve) at (-20,-37) {\faCheckCircle\ Aprovar};
\node[outwarn]  (review)  at (-3,-37)  {\faSearch\ Analisar};
\draw[edge] (income)  -- (history) node[lbl,pos=.55,sloped]{sim};
\draw[edge] (income)  -- (reject)  node[lbl,pos=.5,sloped]{não};
\draw[edge] (history) -- (approve) node[lbl,pos=.55,sloped]{sim};
\draw[edge] (history) -- (review)  node[lbl,pos=.5,sloped]{não};
\end{tikzpicture}}}

\begin{document}

\begin{frame}\titlepage\end{frame}

\begin{frame}{Roteiro da Aula}
  \tableofcontents[hideallsubsections]
\end{frame}

\section{Introdução às Árvores}
\subsection{Objetivos de Aprendizagem}
\begin{frame}{Objetivos de Aprendizagem}
  \begin{itemize}
    \item O que é uma árvore e como ela difere de estruturas lineares?
    \item O que caracteriza uma árvore binária e uma árvore binária de busca?
    \item Quando uma ABB é eficiente? Quais são seus percursos principais?
    \item Por que a AVL precisa de rotações? Onde árvores de decisão se aplicam?
  \end{itemize}
\end{frame}

\subsection{Conceitos Fundamentais}
\begin{frame}{Por que usar árvores?}
  \begin{columns}[T]
    \column{.55\textwidth}
    \begin{itemize}
      \item Listas, pilhas e filas são estruturas \alert{lineares}: guardam os dados em uma única sequência.
      \item Diretórios, organogramas e árvores genealógicas têm relações hierárquicas, não lineares.
      \item Uma árvore representa essas hierarquias diretamente: um mesmo item pode ter vários outros abaixo dele.
    \end{itemize}
    \column{.4\textwidth}\centering
    \FileTreeDiagram
  \end{columns}
\end{frame}

\begin{frame}{Definição formal de uma árvore}
  \begin{itemize}
    \item Uma árvore é vazia ou consiste em uma raiz e zero ou mais subárvores disjuntas.
    \item Cada raiz de subárvore conecta-se à raiz principal por uma aresta.
    \item Formalmente: $T=(r,S)$, em que $r$ é a raiz e $S=\{T_1,T_2,\ldots,T_n\}$ é o conjunto de subárvores.
  \end{itemize}
  \vfill\centering
  \ABC
\end{frame}

\subsection{Nomenclatura Essencial}
\begin{frame}{Nó raiz, pai, filho e irmãos}
  \begin{columns}[T]
    \column{.42\textwidth}
    \begin{itemize}
      \item \textbf{Raiz}: único nó sem pai (A).
      \item \textbf{Pai e filho}: B é pai de D e E.
      \item \textbf{Irmãos}: D e E compartilham o mesmo pai, B.
      \item \textbf{Nível}: distância até a raiz; A está no nível 0.
    \end{itemize}
    \column{.54\textwidth}\centering
    \ABCNiveis
  \end{columns}
\end{frame}

\begin{frame}{Folhas, nós internos, ancestral e descendente}
  \begin{columns}
    \column{.48\textwidth}
    \begin{itemize}
      \item \textbf{Folha}: não possui filhos (D, E, F, G).
      \item \textbf{Nó interno}: possui ao menos um filho (A, B, C).
      \item $u$ é \textbf{ancestral} de $v$ se pertence ao caminho da raiz até $v$; $v$ é descendente de $u$.
    \end{itemize}
    \column{.47\textwidth}\centering
    \ABC\\[1mm]
    \scriptsize A é ancestral de todos os outros nós.
  \end{columns}
\end{frame}

\begin{frame}{Grau, nível e altura}
  \begin{columns}
    \column{.48\textwidth}
    \begin{itemize}
      \item \textbf{Grau}: número de filhos; o grau de B é 2.
      \item \textbf{Nível}: número de arestas desde a raiz.
      \item \textbf{Altura}: maior nível entre as folhas; aqui, $h=2$.
    \end{itemize}
    \column{.47\textwidth}\centering
    \ABC
  \end{columns}
\end{frame}

\section{Árvore Binária}
\subsection{Conceito e Estrutura}
\begin{frame}{Árvore binária}
  \begin{itemize}
    \item Cada nó interno possui \alert{no máximo dois filhos}.
    \item Eles são distintos: subárvore esquerda (LST) e direita (RST).
    \item Formalmente, $T=(x,L,R)$, em que $x$ é a raiz e $L,R$ são árvores binárias disjuntas.
  \end{itemize}
  \vfill\centering
  \ABC\\
  \scriptsize A ordem esquerda--direita importa.
\end{frame}

\subsection{Tipos de Árvores Binárias}
\begin{frame}{Árvore estritamente binária}
  \begin{columns}
    \column{.53\textwidth}
    \textbf{Definição}\par\medskip
    Todo nó possui 0 filhos (folha) ou exatamente 2 filhos; nenhum possui apenas 1 filho.\par\medskip
    Se há $f$ folhas, há $2f-1$ nós.
    \column{.4\textwidth}\centering
    \begin{tikzpicture}[arv,baseline=(current bounding box.center)]
      \node[node]{A}
        child{node[leaf]{B}}
        child{node[node]{C} child{node[leaf]{F}} child{node[leaf]{G}}};
    \end{tikzpicture}
  \end{columns}
\end{frame}

\begin{frame}{Árvore binária cheia (perfeita)}
  \begin{columns}
    \column{.51\textwidth}
    \textbf{Definição}\par\medskip
    Todo nó interno possui dois filhos e todas as folhas estão no mesmo nível. Se a altura é $h$, há $2^{h+1}-1$ nós.
    \column{.44\textwidth}\centering
    \ABC
  \end{columns}
\end{frame}

\begin{frame}{Árvore binária completa}
  \begin{columns}
    \column{.53\textwidth}
    \textbf{Definição}\par\medskip
    Todos os níveis, exceto talvez o último, estão preenchidos; no último, os nós ficam o mais à esquerda possível. Essa forma sustenta os heaps binários.
    \column{.4\textwidth}\centering
    \begin{tikzpicture}[arv,baseline=(current bounding box.center)]
      \node[node]{A}
        child{node[node]{B} child{node[leaf]{D}} child{node[leaf]{E}}}
        child{node[node]{C} child{node[leaf]{F}} child[missing]{}};
    \end{tikzpicture}
  \end{columns}
\end{frame}

\subsection{Árvore Binária de Busca (ABB)}
\begin{frame}{Árvore binária de busca (ABB)}
  \begin{columns}
    \column{.5\textwidth}
    \begin{itemize}
      \item Na subárvore esquerda de $x$, os valores são menores que $x$.
      \item Na direita, os valores são maiores ou iguais a $x$ (convenção adotada aqui).
      \item A propriedade de ordem indica qual ramo explorar em cada passo.
    \end{itemize}
    \column{.45\textwidth}\centering
    \BST
  \end{columns}
\end{frame}

\subsection{Operações em ABB}
\begin{frame}{Operação de inserção em ABB}
  \begin{itemize}
    \item A inserção começa na raiz; o novo nó sempre termina como folha.
    \item Compare a chave: menor segue à esquerda, maior ou igual segue à direita, até chegar a um ponteiro nulo.
    \item Custo médio em árvore balanceada: $O(\log n)$; no pior caso, degenerado: $O(n)$.
  \end{itemize}
\end{frame}

\begin{frame}{Inserção do nó 45: passos 1 e 2}
  \begin{columns}
    \column{.5\textwidth}\centering
    \begin{tikzpicture}[arv,baseline=(current bounding box.center)]
      \node[key]{50}
        child{node[key]{30} child{node[leaf]{20}} child{node[key]{40}}}
        child{node[node]{70}};
    \end{tikzpicture}
    \column{.45\textwidth}
    \begin{enumerate}
      \item $45<50$: desça à esquerda.
      \item $45>30$: desça à direita.
    \end{enumerate}
  \end{columns}
\end{frame}

\begin{frame}{Inserção do nó 45: passos 3 e 4}
  \begin{columns}
    \column{.5\textwidth}\centering
    \begin{tikzpicture}[arv,baseline=(current bounding box.center)]
      \node[key]{50}
        child{node[key]{30} child{node[leaf]{20}}
          child{node[key]{40} child[missing]{} child{node[leaf,draw=treeorange,fill=orange!18]{45}}}}
        child{node[node]{70}};
    \end{tikzpicture}
    \column{.45\textwidth}
    \begin{enumerate}
      \setcounter{enumi}{2}
      \item $45>40$: tente ir à direita.
      \item A posição é nula: insira 45.
    \end{enumerate}
  \end{columns}
\end{frame}

\begin{frame}{Operação de remoção em ABB}
  \begin{enumerate}
    \item \textbf{Folha}: desconecte o nó.
    \item \textbf{Um filho}: conecte o pai diretamente ao filho do nó removido.
    \item \textbf{Dois filhos}: use o sucessor em ordem (menor da direita) ou o predecessor (maior da esquerda); copie a chave e remova-o do local original.
  \end{enumerate}
\end{frame}

\begin{frame}{Remoção: caso 1 (nó folha)}
  \centering
  \DeleteLeafDiagram
\end{frame}

\begin{frame}{Remoção: caso 2 (um filho)}
  \centering
  \DeleteOneChildDiagram
\end{frame}

\begin{frame}{Remoção: caso 3 (dois filhos)}
  \begin{columns}
    \column{.44\textwidth}\small
    Ao remover 15, o sucessor em ordem é 20, o menor valor da subárvore direita. Copie essa chave para a raiz e remova a folha 20 original.
    \column{.51\textwidth}\centering
    \DeleteTwoChildrenDiagram
  \end{columns}
\end{frame}

\subsection{Algoritmos de Percurso}
\begin{frame}{Percursos em árvores binárias}
  Todos visitam cada nó uma vez: custo $O(n)$. O que muda entre eles é apenas a posição da visita à raiz:
  \begin{itemize}
    \item Pré-ordem: raiz, esquerda, direita.
    \item Em ordem: esquerda, raiz, direita.
    \item Pós-ordem: esquerda, direita, raiz.
  \end{itemize}
\end{frame}

\begin{frame}{Percurso em pré-ordem (Raiz, Esquerda, Direita)}
  \begin{columns}
    \column{.44\textwidth}
    \begin{enumerate}
      \item Visite a raiz.
      \item Percorra a esquerda.
      \item Percorra a direita.
    \end{enumerate}
    \alert{A, B, D, E, C, F, G}.
    \column{.5\textwidth}\centering
    \ABC
  \end{columns}
\end{frame}

\begin{frame}{Percurso em ordem (Esquerda, Raiz, Direita)}
  \begin{columns}
    \column{.44\textwidth}
    \begin{enumerate}
      \item Percorra a esquerda.
      \item Visite a raiz.
      \item Percorra a direita.
    \end{enumerate}
    \alert{D, B, E, A, F, C, G}. Em uma ABB, as chaves saem em ordem crescente.
    \column{.5\textwidth}\centering
    \ABC
  \end{columns}
\end{frame}

\begin{frame}{Percurso em pós-ordem (Esquerda, Direita, Raiz)}
  \begin{columns}
    \column{.44\textwidth}
    \begin{enumerate}
      \item Percorra a esquerda.
      \item Percorra a direita.
      \item Visite a raiz.
    \end{enumerate}
    \alert{D, E, B, F, G, C, A}.
    \column{.5\textwidth}\centering
    \ABC
  \end{columns}
\end{frame}

\subsection{Implementação}
\begin{frame}[fragile]{Implementação de uma ABB em Python}
\begin{lstlisting}[language=Python]
class No:
    def __init__(self, chave):
        self.chave, self.esq, self.dir = chave, None, None

def inserir(raiz, chave):
    if raiz is None: return No(chave)
    if chave < raiz.chave:
        raiz.esq = inserir(raiz.esq, chave)
    else: raiz.dir = inserir(raiz.dir, chave)
    return raiz
\end{lstlisting}
\end{frame}

\begin{frame}[fragile]{Percurso em ordem}
\begin{lstlisting}[language=Python]
def em_ordem(no):
    if no is not None:
        em_ordem(no.esq)
        print(no.chave, end=" ")
        em_ordem(no.dir)
\end{lstlisting}
Em uma ABB, essa função imprime as chaves em ordem não decrescente.
\end{frame}

\section{Árvore AVL}
\subsection{O Problema do Desbalanceamento}
\begin{frame}{O problema do desbalanceamento}
  \begin{columns}
    \column{.48\textwidth}\centering
    \textbf{Balanceada}\\\scriptsize Busca $O(\log n)$\\[1mm]
    \begin{tikzpicture}[
        level distance=9mm,level 1/.style={sibling distance=16mm},level 2/.style={sibling distance=7mm},
        edge from parent/.style={draw=treeblue,thick}]
      \node[node]{4}
        child{node[node]{2} child{node[leaf]{1}} child{node[leaf]{3}}}
        child{node[node]{6} child{node[leaf]{5}} child{node[leaf]{7}}};
    \end{tikzpicture}
    \column{.48\textwidth}\centering
    \textbf{Degenerada}\\\scriptsize Busca $O(n)$\\[1mm]
    \begin{tikzpicture}[level distance=8mm,sibling distance=10mm,edge from parent/.style={draw=treeblue,thick}]
      \node[node]{1}
        child[missing]{}
        child{node[node]{2} child[missing]{} child{node[node]{3} child[missing]{} child{node[leaf]{4}}}};
    \end{tikzpicture}
  \end{columns}
  \vfill\centering
  \alert{Solução: ABB autobalanceável, como a AVL.}
\end{frame}

\subsection{Conceito de Árvore AVL}
\begin{frame}{Árvore AVL: definição}
  \begin{itemize}
    \item Uma AVL (Adelson-Velskii e Landis) é uma ABB com uma condição extra sobre a altura das subárvores.
    \item $FB(no)=h(\text{direita})-h(\text{esquerda})$.
    \item Em uma AVL, $FB\in\{-1,0,1\}$; os valores $-2$ ou $2$ exigem rebalanceamento.
    \item Com essa condição, busca, inserção e remoção ficam em $O(\log n)$.
  \end{itemize}
\end{frame}

\subsection{Rotações em Árvores AVL}
\begin{frame}{Rebalanceamento com rotações}
  Uma rotação reorganiza poucos nós, preserva a ordem da ABB e recupera o equilíbrio de altura.
  \medskip
  \centering\small
  \begin{tabular}{lll}
    \toprule
    \textbf{FB do nó} & \textbf{FB do filho pesado} & \textbf{Rotação} \\
    \midrule
    $+2$ & $0$ ou $+1$ & simples à esquerda (caso RR) \\
    $+2$ & $-1$        & dupla direita--esquerda (caso RL) \\
    $-2$ & $0$ ou $-1$ & simples à direita (caso LL) \\
    $-2$ & $+1$        & dupla esquerda--direita (caso LR) \\
    \bottomrule
  \end{tabular}
\end{frame}

\begin{frame}{Rotação simples à esquerda}
  \begin{columns}
    \column{.46\textwidth}\small
    Após inserir 12, o nó 6 fica com $FB(6)=+2$. Como o caminho pesado segue para a direita de 8, temos o caso RR: a rotação coloca 8 na raiz local.
    \column{.49\textwidth}\centering
    \RotationLeftDiagram
  \end{columns}
\end{frame}

\begin{frame}{Rotação dupla direita--esquerda}
  \begin{columns}
    \column{.46\textwidth}\small
    Ao inserir 7, surge um caminho em "joelho": $FB(6)=+2$ e $FB(8)=-1$. Primeiro rotacionamos 8 à direita; em seguida, 6 à esquerda.
    \column{.49\textwidth}\centering
    \RotationRLDiagram
  \end{columns}
\end{frame}

\section{Árvores de Decisão: aplicação atual}
\subsection{Conceito, avaliação e estado da arte}
\begin{frame}{Árvore de decisão: da estrutura ao modelo}
  \begin{columns}
    \column{.52\textwidth}
    \begin{itemize}
      \item No aprendizado supervisionado, cada nó testa uma variável; as arestas registram os resultados possíveis, e as folhas retornam uma classe ou um valor.
      \item O treinamento procura divisões que reduzam a impureza (entropia ou índice Gini) ou o erro de regressão.
      \item Profundidade limitada, poda e validação em dados separados ajudam a conter o sobreajuste.
    \end{itemize}
    \column{.43\textwidth}\centering
    \DecisionTreeDiagram
  \end{columns}
\end{frame}

\begin{frame}{Avaliação responsável}
  \begin{itemize}
    \item Separe treino, validação e teste. Ajuste hiperparâmetros somente com os dados de validação.
    \item Quando as classes são desbalanceadas, acompanhe precisão, revocação, F1, PR-AUC e a matriz de confusão -- acurácia isolada pode esconder erros relevantes.
    \item Procure vazamento de alvo e compare o desempenho entre subgrupos que façam sentido para o problema.
    \item Decisões de alto impacto exigem revisão humana, registro das escolhas e acompanhamento após a implantação.
  \end{itemize}
\end{frame}

\begin{frame}{Árvores em métodos atuais para dados tabulares}
  \small
  \begin{itemize}
    \item Uma árvore isolada é fácil de explicar. Florestas aleatórias combinam muitas árvores para reduzir a variância.
    \item Em \textit{gradient boosting}, cada árvore nova corrige os erros do conjunto anterior; XGBoost, LightGBM e CatBoost são implementações amplamente usadas.
    \item Esses conjuntos costumam ganhar em desempenho, mas são mais difíceis de explicar do que uma árvore isolada -- a comparação deve seguir o protocolo de avaliação do problema.
  \end{itemize}
  \vfill\scriptsize Fontes: Chen e Guestrin (2016); Ke et al. (2017); Prokhorenkova et al. (2018).
\end{frame}

\section{Conclusões}
\begin{frame}{Conclusões}
  \begin{itemize}
    \item Árvores modelam hierarquias; árvores binárias limitam cada nó a dois filhos ordenados.
    \item ABBs são eficientes quando a altura é pequena; a AVL garante $O(\log n)$ com rotações.
    \item Pré-ordem, em ordem e pós-ordem custam $O(n)$ e servem a objetivos distintos.
    \item Árvores de decisão aplicam essa estrutura a classificação e regressão; avaliação e governança são essenciais.
  \end{itemize}
\end{frame}

\section{Referências Bibliográficas}
\begin{frame}[allowframebreaks]{Referências}
\scriptsize
\begin{thebibliography}{99}
\bibitem{lee} Lee, K. D.; Hubbard, S. \textit{Data Structures and Algorithms with Python}. Springer, 2015.
\bibitem{clrs} Cormen, T. H. et al. \textit{Introduction to Algorithms}. 3. ed. MIT Press, 2009.
\bibitem{sz} Szwarcfiter, J. L.; Markenzon, L. \textit{Estruturas de Dados e seus Algoritmos}. LTC, 2015.
\bibitem{breiman} Breiman, L. Random Forests. \textit{Machine Learning}, 45, 5--32, 2001.
\bibitem{xgb} Chen, T.; Guestrin, C. XGBoost: A Scalable Tree Boosting System. \textit{KDD}, 2016. \url{https://doi.org/10.1145/2939672.2939785}
\bibitem{lgb} Ke, G. et al. LightGBM: A Highly Efficient Gradient Boosting Decision Tree. \textit{NeurIPS}, 2017. \url{https://proceedings.neurips.cc/paper/2017/hash/6449f44a102fde848669bdd9eb6b76fa-Abstract.html}
\bibitem{cat} Prokhorenkova, L. et al. CatBoost: unbiased boosting with categorical features. \textit{NeurIPS}, 2018. \url{https://proceedings.neurips.cc/paper/2018/hash/14491b756b3a51daac41c24863285549-Abstract.html}
\end{thebibliography}
\end{frame}

\begin{frame}\titlepage\end{frame}

\end{document}
